\centerline{\bf \Huge Мега теория по многочленам}

\begin{flushright}
    \scriptsize{Я ждал этого момента 4...\\нет 5 тысяч лет!}
\end{flushright}

\textbf{\large \bf Персонаж номер 1. Дело кольца целых чисел.}

Прежде всего давайте вспомним некоторые определения и факты про обычные целые числа, чтобы мы потом могли проводить аналогии с многочленами.

Целое число $a$ называется \textit{обратимым}, если существует такое целое число $b$, что $ab=1$.

Нетрудно понять, что обратимые элементы в $\mathbb{Z}$ это $\pm 1$.

Целое число называется \textit{простым}, если оно делится только на себя или на \textbf{обратимый элемент}(!).

\textbf{Основная теорема арифметики для $\mathbb{Z}$.}

Любое целое число $a$ представляется в виде произведения степеней простых чисел с точностью до перестановки сомножителей и домножения на обратимые элементы. Иначе говоря, $a = p_1^{\alpha_1}\cdot p_2^{\alpha_2} \dotsc \cdot p_k^{\alpha_k}$. 

Поняв с помощью ОТА структуры целых чисел, мы сразу же вводим такие понятия, как НОД и НОК и первое оказывается наиболее интересным. Для целых чисел существует так называемый алгоритм Евклида, помогающий найти НОД двух чисел. Вспомним же, как он работает и заодно вспомним, что такое линейная представимость НОД.

\textbf{Алгоритм Евклида для $\mathbb{Z}$}

Докажем, что $\text{НОД}(a,b)=\text{НОД}(a-b,b)$.

Покажем, что все общие делители чисел $a$ и $b$ и чисел $a-b$ и $b$ совпадают: пусть $d|a$ и $d|b$, тогда $d|(a-b)$. Также, если $d|(a-b)$ и $d|b$, тогда $d|((a-b)+b)$.

Из этих утверждений получается, что множества общих делителей чисел $a,b$ и чисел $(a-b),b$ совпадают, значит, совпадают и их наибольшие элементы. 

Повторив вычитание много раз, дойдем до остатка от деления $a$ на $b$. Получается, что  $\text{НОД}(a,b)=\text{НОД}(r,b)$, где $r$ - остаток от деления числа $a$ на $b$.

Алгоритм Евклида заключается в том, чтобы повторить операцию описанную выше много раз: 

$(a,b)=(r_1,b)=(r_1,r_2)=\ldots=(r_{n-1},r_n)$.

\newpage

\hspace{3mm}

Где

$a=b\cdot q_1+r_1$

$b=r_1 \cdot q_2 + r_2$

$r_1=r_2 \cdot q_3 + r_3$

$\ldots$

$r_{n-1}=r_n \cdot q_{n+1} $

Процесс сходится, так как $|r_i| > |r_{i+1}|$, и число $|b|$ ограничено, бесконечный спуск невозможен.

\textbf{Линейная представимость НОД для $\mathbb{Z}$}

Для любых целых $a$,$b$ существуют такие числа $x$ и $y$, что $ax+by=\text{НОД}(x,y)$.

Распишем алгоритм Евклида:

$(a,b)=(r_1,b)=(r_1,r_2)=\ldots=(r_{n-1},r_n)$, где 

$a=b\cdot q_1+r_1$

$b=r_1 \cdot q_2 + r_2$

$r_1=r_2 \cdot q_3 + r_3$

$\ldots$

$r_{n-1}=r_n \cdot q_{n+1} $


Выразим из первого уравнения $r_1=a-b\cdot q_1$ и подставим во второе:

$b=(a-b \cdot q_1) q_2 + r_2$. Отсюда выразим $r_2=b-(a-b\cdot q_1) \cdot q_2$.

Подставим все в третье уравнение:

$(a-b\cdot q_1)=(b-(a-b\cdot q_1) \cdot q_2) \cdot q_3 + r_3$ откуда мы можем выразить $r_3$ и так далее, вплоть до $r_n$. 


{\bf Пример}

Найдем НОД (28, 77).

$77=28\cdot2 +21$

$28=21 \cdot 1 +7$

$21=7 \cdot 3$

Выразим из первого $21=77-28\cdot2$

$28=(77-28\cdot2)\cdot 1+7$ откуда $28-(77-28\cdot 2)=7$

Отсюда получаем, что $7=28\cdot 3 - 77$.

Это и есть линейная представимость 7 через 28 и 77. 

$7=28\cdot 3 - 77$

\newpage

\hspace{3mm}

\textbf{\large \bf Персонаж номер 2. Дело многочленов над целыми числами}

Теперь, вспомнив основные понятия из целых чисел можем переходить к многочленам. До этого мы рассматривали множество $\mathbb{Z}$ - всех целых чисел. Теперь будем рассматривать множество $\mathbb{Z}[x]$ всех многочленов с целыми коэффициентами от переменной $x$. Иначе это множество называются \textit{многочлены над кольцом целых чисел.}

Как мы знаем константы(то есть сами целые числа) это тоже многочлены, просто степени 0. Такие многочлены мы будем называть \textit{тривиальными}.

Теперь приведем аналог простых чисел среди многочленов.

Многочлен $P(x)\neq 0$ называется \textit{неприводимым}, если он не раскладывается в произведение двух нетривиальных многочленов $P(x) = Q(x)K(x)$.

\textit{Упражнение.} Приведите примеры неприводимых многочленов над целыми числами.

\textbf{Основная теорема арифметики для $\mathbb{Z}[x]$}

Любой многочлен над кольцом целых чисел представляется в виде произведение неприводимых многочленов с точностью до перестановки сомножителей и домножения на обратимые элементы (то есть на $\pm 1$).

\textbf{Наибольшой общим делителем} многочленов $P(x)$ и $Q(x)$ называется многочлен наибольшей степени $D(x)$ такой, что $P(x)$ и $Q(x)$ делятся на $D(x)$. Обозначают точно так же, как и в целых числах $D(x) = (P(x), Q(x))$. НОД двух многочленов определен точно так же с точностью до домножения на $\pm 1$

НОД можно найти точно так же с помощью алгоритма Евклида.

\textbf{Алгоритм Евклида для $\mathbb{Z}[x]$}


Мы уже хорошо умеем делить многочлены с остатком и хорошо понимаем, что НОД$(P(x), Q(x)) = $ НОД$(P(x) - Q(x), Q(x))$. Поэтому давайте сразу делить многочлен $P(x)$ на $Q(x)$ с остатком и этот самый остаток писать вместо $P(x)$ точно так же, как мы делали с целыми числами.

\centerline{$P(x) = Q(x)K(x)+R_0(x) -$ нулевой шаг}

\centerline{$K(x) = K_1(x)R_0(x) + R_1(x) -$ первый шаг}

\centerline{$\dotsc$}

\centerline{$R_{n-2}(x) = K_n(x)R_{n-1}(x) + R_n - n-1$-й шаг}

\centerline{$R_{n-1}(x) = K_{n+1}(x)R_{n}(x) - n$-й шаг}

Здесь $R_n(x)$ это последний ненулевой остаток, который и будет являться НОДом наших многочленов.

\textit{Упражнение.} Найдите НОД$(x^6+3x^5+2x+1, x^2+1)$

\hspace{3mm}

Напоследок сформулируем теорему о линейном представлении НОД двух многочленов.

\textbf{Линейное представление НОД для $\mathbb{Z}[x]$}

Пусть $D(x) = (P(x), Q(x)),$ где $P(x), Q(x) \in Z[x]$, тогда существует многочлены $U(x), V(x)$ такие что: 

$D(x) = U(x)P(x)+V(x)Q(x)$